\section{Related work}
\label{sec:related}
\paragraph{Reconstruction and perceptual quality.}
The perception--distortion tradeoff separates reference distortion from distributional perceptual quality~\cite{blau2018}.
Learned Perceptual Image Patch Similarity (LPIPS) evaluates aspects of visual similarity that pixel distances miss~\cite{zhang2018lpips}.
Neither objective alone defines whether the retained detail belongs to the target person.
Our audit adds controlled person-, organ-, and acquisition-level comparisons to this distinction.
Structural similarity (SSIM) supplies a local reference comparison~\cite{wang2004}, but its interpretation still depends on the anatomical and acquisition changes represented by a pair.

\paragraph{Anatomical completeness and harmonization.}
CARE evaluates sparse-view CT reconstruction using anatomy-aware measurements of large and small organs, intestines, and vessels~\cite{care2025}.
Its objective is structural completeness rather than validation of a donor-based identity scale.
LMSB addresses anatomical changes in OCT harmonization by learning an anatomy-guided latent metric before distribution matching~\cite{lmsb2025}.
It demonstrates why Euclidean proximity in appearance need not coincide with anatomical consistency.
Our contribution is an evaluation design around such interpretations, not an empirical comparison against either reconstruction or harmonization system.

\paragraph{Anatomical representations and biometric units.}
Brain-ID learns contrast-agnostic anatomical features from synthetic inputs and evaluates robustness and downstream brain-imaging tasks~\cite{brainid2024}.
Those results motivate acquisition controls, but a Brain-ID distance has not been completed in this study.
Iris recognition illustrates a complementary distinction: a person has two eyes, and person membership does not make their iris textures one biometric instance~\cite{daugman2004,daugman2001}.
The relevant identity unit must therefore be specified before an anatomical similarity result is interpreted.
ArcFace supplies an identity-oriented face representation~\cite{arcface2022}; our local-donor analysis distinguishes that whole-image reference from the window-level reconstruction distance.

\paragraph{Task-specific metric validation.}
Metrics Reloaded emphasizes matching validation measures to the task and data~\cite{metrics2024}.
We operationalize that principle through separate controls for attenuation, incorrect detail, bilateral correspondence, and acquisition.
A successful check on one axis does not establish success on the others.
